\documentclass[11pt]{article}

\usepackage[margin=0.82in]{geometry}
\usepackage[english]{babel}
\usepackage[utf8]{inputenc}
\usepackage[T1]{fontenc}
\usepackage{lmodern}
\usepackage{amsmath,amssymb,bm}
\usepackage{booktabs}
\usepackage{enumitem}
\usepackage{xcolor}
\usepackage{tabularx}
\usepackage{longtable}
\usepackage{array}
\usepackage{fancyhdr}
\usepackage{microtype}
\usepackage{listings}
\usepackage{xspace}
\usepackage[
  colorlinks=true,
  linkcolor=blue!55!black,
  urlcolor=blue!55!black,
  citecolor=blue!55!black
]{hyperref}

\definecolor{uclablue}{HTML}{2774AE}
\definecolor{codegray}{HTML}{F5F6F7}
\definecolor{darkgray}{HTML}{3F4850}

\pagestyle{fancy}
\fancyhf{}
\lhead{\textit{ECE 232E --- Summer 2026}}
\rhead{\textit{PANINI Course Project}}
\cfoot{\thepage}
\setlength{\headheight}{14pt}

\setlist[itemize]{leftmargin=1.5em,itemsep=2pt,topsep=3pt}
\setlist[enumerate]{leftmargin=1.8em,itemsep=3pt,topsep=3pt}
\setlist[description]{leftmargin=1.4em,itemsep=3pt,topsep=3pt}

\newcolumntype{Y}{>{\raggedright\arraybackslash}X}
\newcolumntype{C}{>{\centering\arraybackslash}p{0.13\textwidth}}
\newcommand{\code}[1]{\texttt{#1}}
\newcommand{\RICR}{{\mdseries\textsc{RICR}}\xspace}
\newcommand{\PANINI}{{\mdseries\textsc{Panini}}\xspace}
\newcommand{\question}[3]{%
  \subsection[Question #1: #2 (#3 points)]{Question #1: #2 \hfill [#3 points]}%
}
\newcommand{\deliverables}{\paragraph{Required output.}}
\newcommand{\grading}{\paragraph{Point allocation.}}

\lstset{
  basicstyle=\ttfamily\small,
  backgroundcolor=\color{codegray},
  frame=single,
  rulecolor=\color{black!18},
  breaklines=true,
  columns=fullflexible,
  showstringspaces=false,
  keywordstyle=\color{uclablue}\bfseries,
  commentstyle=\color{darkgray}
}

\begin{document}

\begin{center}
{\large\bfseries Large Scale Social and Complex Networks: Design and Algorithms}\\[2pt]
{\large\bfseries ECE 232E --- Summer 2026}\\[3pt]
Prof.\ Vwani Roychowdhury \hfill UCLA Department of Electrical and Computer Engineering
\end{center}
\vspace{2mm}
\hrule height 0.8pt
\vspace{4mm}
\begin{center}
  {\Huge\bfseries Course Project}\\[3pt]
  {\LARGE\bfseries Structured Memory Networks and}\\[2pt]
  {\LARGE\bfseries Reasoning Inference Chain Retrieval}\\[5pt]
  {\large Based on \PANINI: Continual Learning in Token Space via Structured Memory}\\[5pt]
  {\large\bfseries 100 points}\\[2pt]
  {\large Due Saturday, August 15, 2026 at 11:59 PM Pacific Time}
\end{center}
\vspace{4mm}
\hrule height 0.8pt
\vspace{5mm}

\section{Purpose and learning objectives}

Standard retrieval-augmented generation stores document chunks and asks a
language model to repeatedly reason over retrieved text. \PANINI instead
converts documents into GSWs: heterogeneous semantic networks containing
entities, verb-phrase/event nodes, roles, states, and grounded question--answer
(QA) edges. At inference time, \RICR decomposes a multi-hop query and searches
for compact chains of QA pairs through this structured memory.

This project connects the \PANINI framework to core topics from ECE 232E:
heterogeneous graph construction, network projection, connectivity,
centrality, ranking, sparse and dense search, beam search, scaling, and
controlled experimental evaluation. By the end of the project, you should be
able to:

\begin{enumerate}
  \item construct and characterize corpus-level networks from document-level
        GSWs;
  \item compare TF--IDF, BM25, dense, hybrid, and structure-aware retrieval;
  \item execute and validate a neural question-decomposition plan;
  \item implement \RICR with chain scoring and diversity-preserving pruning;
  \item analyze generalization from 2WikiMultiHopQA to MuSiQue; and
  \item report accuracy, retrieval quality, latency, evidence size, and
        failure modes reproducibly.
\end{enumerate}

\section{Reading, repository, and compute}

\paragraph{Required reading.}
Read Sections 1--4 and Algorithm 1 of
\href{https://arxiv.org/abs/2602.15156}{\PANINI: Continual Learning in Token
Space via Structured Memory}. The GSW representation, dual indexing,
decomposition notation, chain scoring, and \RICR pruning rule in the paper are
part of the project specification.

\paragraph{Student repository.}
The self-contained starter repository is:
\begin{center}
\url{https://github.com/YigitTurali/panini-course-project}
\end{center}
The Colab notebook, artifact loaders, graph utilities, metrics, embeddings,
indices, and model configuration are supplied. Students do \textbf{not} need
access to the research codebase.

\paragraph{Models.}
Use the model identifiers in \code{models/model\_config.json}. The required
free-tier configuration uses:
\begin{itemize}
  \item \href{https://huggingface.co/yigitturali/GSW-QA-Decomposer-Qwen3-4B}
        {GSW-QA-Decomposer-Qwen3-4B} for decomposition;
  \item supplied Qwen3-Embedding-8B corpus and fixed-query embeddings;
  \item \href{https://huggingface.co/Qwen/Qwen3-Reranker-8B}
        {Qwen3-Reranker-8B} in 4-bit mode for reranking; and
  \item \href{https://huggingface.co/Qwen/Qwen3-4B}
        {Qwen3-4B} in 4-bit mode for evidence-grounded final answers.
\end{itemize}
The
\href{https://huggingface.co/yigitturali/GSW-QA-Decomposer-Qwen3-8B}
{8B decomposer} is optional and receives no automatic accuracy advantage.
Do not keep the decomposer, reranker, query encoder, and answer model resident
on the GPU simultaneously. Save outputs to JSONL, unload the current model,
clear GPU memory, and then begin the next stage.

\paragraph{Compute policy.}
The required path must run on a free Google Colab GPU. No paid API is required
or permitted for the required results. Corpus GSWs, embeddings, and indices
are supplied; regenerating them is outside the project scope.

\section{Datasets and evaluation protocol}

You will receive two independent 100-question packages with the same directory
contract and stable IDs.

\begin{table}[h]
\centering
\small
\begin{tabularx}{\textwidth}{@{}l C C C Y@{}}
\toprule
\textbf{Dataset} & \textbf{Development} & \textbf{Held out} & \textbf{Total}
& \textbf{Required stratification}\\
\midrule
2WikiMultiHopQA & 80 & 20 & 100 &
25 bridge-comparison, 25 comparison, 25 compositional, and 25 inference
questions; each category has 20 development and 5 held-out questions.\\
MuSiQue & 80 & 20 & 100 &
50 two-hop, 30 three-hop, and 20 four-hop questions; development/held-out
counts are respectively 40/10, 24/6, and 16/4.\\
\bottomrule
\end{tabularx}
\caption{Required course subsets.}
\end{table}

\begin{itemize}
  \item The \textbf{development split} includes questions, answers, aliases,
        supporting evidence, and context-document IDs. It may be used for
        debugging and tables in the report.
  \item The \textbf{held-out split} includes questions and document IDs but
        withholds answers and supporting evidence. Do not manually search for,
        reconstruct, or import these labels.
  \item Select algorithms and default hyperparameters using only the 2Wiki
        development split. Freeze them before running MuSiQue. MuSiQue is the
        required cross-dataset transfer evaluation.
  \item Run the final frozen system once on all 200 questions. Report metrics
        locally for the 160 development questions; the teaching staff will
        score the 40 held-out predictions.
\end{itemize}

MuSiQue is distributed under
\href{https://github.com/stonybrooknlp/musique}{CC BY 4.0}; 2WikiMultiHopQA is
distributed under
\href{https://github.com/Alab-NII/2wikimultihop}{Apache 2.0}. Retain the
dataset IDs and attribution files in all derived outputs.

\section{Fixed definitions and default settings}

\subsection{Graph views and the reconciliation boundary}

For each GSW, construct the \textbf{native directed multigraph}. Entity,
space/time, and verb-phrase/event records are nodes. Each grounded QA answer is
a directed edge from its verb-phrase node to its answer node. The edge retains
the stable QA ID, question, document ID, and direction. Do not merge records
merely because their display strings match.

Next construct an \textbf{unreconciled entity projection}. Its nodes are the
document-local entity occurrences identified by their stable UIDs. Two
occurrences are adjacent when they answer QA records attached to the same
local verb-phrase; repeated neighborhoods increment edge weight. Finally,
construct one or more \textbf{analysis-only reconciled projections} by mapping
occurrence UIDs to proposed cross-document entity IDs and aggregating the
projected edges.

The supplied exact-surface baseline case-folds, removes punctuation, and
collapses spaces. Students must implement and justify a more conservative
reconciliation rule that also uses available type, role/state, and local
neighborhood evidence. Neither method has complete gold identity labels, so
evaluate it with a reproducible manual audit and a sensitivity comparison.
All network results must name the precise graph and reconciliation rule used.

\paragraph{Why the reconciled graph is not used by \PANINI.}
The paper reconciles entities within a document but does not precompute
cross-document entity links. Its sparse entity index retains the originating
document-local node, its dense index retrieves QA pairs directly, and \RICR
creates cross-document chains at read time by inserting a retrieved answer
into the next sub-question. A global surface merge can incorrectly join
homonyms, titles, years, occupations, and nationality values into artificial
hubs; a conservative rule can instead miss aliases. These errors directly
distort components and centralities and would contaminate retrieval if treated
as facts. In this project reconciliation is therefore a network-analysis
exercise only. Do not use reconciled edges for entity expansion, candidate
generation, reranking, \RICR, or answer prediction.

\subsection{Retrieval and \RICR defaults}

Unless a question explicitly requests an ablation, use:

\begin{center}
\begin{tabular}{@{}ll@{}}
\toprule
Random seed & 232\\
Dense similarity & inner product of L2-normalized vectors\\
Reciprocal Rank Fusion & $\operatorname{RRF}(d)=\sum_r 1/(60+r(d))$\\
Pre-rerank candidate pool & $M=60$ unique QA records\\
Candidates returned per hop & $k=15$\\
Beam width & $B=5$\\
Chain score & geometric mean of hop scores\\
Generation & deterministic; sampling disabled\\
\bottomrule
\end{tabular}
\end{center}

For a length-$t$ chain $C=(e_1,\ldots,e_t)$ with positive hop scores
$s_1,\ldots,s_t$, use
\begin{equation}
  \operatorname{score}(C)
  = \left(\prod_{\ell=1}^{t}\max(s_\ell,\epsilon)\right)^{1/t}
  = \exp\left(\frac{1}{t}\sum_{\ell=1}^{t}
    \log\max(s_\ell,\epsilon)\right),
  \qquad \epsilon=10^{-12}.
\end{equation}
The log-domain form is recommended. At every pruning step, sort by decreasing
chain score and use the stable QA-ID sequence as the deterministic tie-breaker.
Retain at most one chain per normalized \emph{current answer}.

\subsection{Metrics}

Use the provided normalization functions. Report:
\begin{description}
  \item[Retrieval:] Recall@$k$ for $k\in\{1,5,10,15\}$, reciprocal rank/MRR,
        supporting-document recall, and supporting-QA recall.
  \item[Chains:] complete-chain recovery (all required supporting hops occur in
        one surviving chain), average surviving chains, and average unique
        current answers.
  \item[Answers:] Exact Match (EM) and token F1, giving each question the
        maximum score over its accepted answer aliases.
  \item[Efficiency:] wall-clock latency per stage, peak GPU memory when a model
        is used, number of reranked candidates, number of returned QA pairs,
        and answer-model input tokens.
\end{description}
Aggregate by macro-average over questions. Include the number of evaluated
questions in every table. Never average already-rounded subgroup values.

\section{Graded questions}

\question{1}{Understand the two data packages}{4}

Imagine that the write-time stage of \PANINI has just finished processing two
collections of documents. You did not run that stage: you inherit its
document-level GSWs, flattened entity and QA tables, embeddings, and indices.
Before treating these files as one memory, you must establish what one row
means and which identifiers remain valid when records from many documents are
combined. A retrieval result that points to the wrong embedding row can look
plausible while silently invalidating every later experiment.

Load both packages with \code{CoursePackage}. Use \code{entity\_uid} and
\code{qa\_uid} as identifiers rather than list positions. Check that every
embedding row has exactly one matching ID, IDs are unique, and held-out files
contain inputs but no grading labels. This audit becomes the trusted boundary
for all later questions.

\begin{enumerate}[label=(\alph*)]
  \item Produce a table containing, for each dataset and split, the number of
        questions, unique documents, GSW files, entities, QA records, and
        decomposition-validation examples.
  \item Show one question, entity, and QA record. Explain the stable ID and
        every field that your implementation consumes.
  \item Add an assertion that held-out records contain no answer-label fields:
        \code{answer}, \code{answer\_aliases},
        \code{supporting\_facts}, or \code{evidences}.
  \item Run \code{pytest -q tests}. Record which tests pass and which RICR
        tests are skipped before you begin the implementation.
\end{enumerate}

\deliverables One dataset table, three short schema examples, the no-leakage
assertion, and the initial test summary in the notebook.

\grading Dataset table (1), schema explanation (1), automated leakage check
(1), and starter-test run (1).

\question{2}{Build the GSW network and reconcile entities}{8}

The verified records from Question 1 are still hundreds of separate local
workspaces. To study them as a network, first preserve exactly what was written
inside each document, and only afterward ask whether two occurrences refer to
the same real-world entity. Mixing those operations would make it impossible
to tell whether a global edge came from a GSW or from your identity rule.

Construct the native directed multigraph by namespacing every node with its
document and GSW filename. Add entity, space/time, and verb-phrase nodes. For
each grounded QA answer, add a verb-phrase-to-answer edge keyed by QA ID and
retain question text and provenance. Then project each local verb neighborhood
onto its answer entities without merging across documents. Reconciliation
must remain a separate occurrence-UID-to-global-ID mapping so that the same
unreconciled graph can be evaluated under several identity hypotheses.

\begin{enumerate}[label=(\alph*)]
  \item Report native-graph node and edge counts by node/edge type and verify
        them against direct counts from the JSON files.
  \item Construct the weighted unreconciled entity projection. Verify that
        occurrence UIDs from different documents remain distinct, even when
        their names match.
  \item Run the exact-surface baseline, then implement a conservative
        reconciliation rule using surface, type/role compatibility, and local
        neighborhood evidence. State the decision rule precisely.
  \item Manually audit a fixed-seed sample of at least 30 proposed
        cross-document merges. Label each correct, incorrect, or uncertain;
        report estimated precision with uncertain cases shown separately.
        Include at least two false-merge and two missed-alias examples.
  \item Add unit tests for edge direction, provenance, projection weight,
        cross-document non-merge, intended merge, and homonym non-merge.
\end{enumerate}

\deliverables Construction and reconciliation code, the audit table, graph
counts under all three identity conditions, and passing tests.

\grading Correct native graph (2), correct unreconciled projection (2),
reconciliation method and audit (3), and tests (1).

\question{3}{Analyze the structured-memory network}{8}

Question 2 produces several candidate ``global memories,'' but network
statistics can make a bad reconciliation look convincing. For example,
merging every occurrence of a nationality, year, occupation, or ambiguous
title can manufacture a giant component and high-centrality hubs that were
never asserted by any document. Your task is therefore not merely to compute
statistics; it is to decide which observed structure belongs to the local GSWs
and which was created by the reconciliation rule.

Perform the following analysis separately for 2Wiki and MuSiQue. Compute every
requested property on the unreconciled projection, exact-surface projection,
and conservative projection; include native-graph measurements as a reference.

\begin{enumerate}[label=(\alph*)]
  \item Compute weakly connected components of the native graph and connected
        components of each entity projection. Report the number of components,
        giant-component size and fraction, isolated nodes, and component-size
        summary (minimum, median, mean, and maximum). Quantify how much each
        reconciliation rule changes these values.
  \item Plot degree probability mass and complementary cumulative degree
        distributions on log--log axes. State whether the heavy tail, if any,
        is consistent across datasets and reconciliation rules; do not claim a
        power law without a fitted test.
  \item For each reconciled projection, report the ten highest nodes under
        weighted degree, PageRank, and betweenness centrality. Audit whether
        hubs represent entities, common attribute values, or reconciliation
        errors. Exact betweenness or a documented fixed-seed approximation is
        acceptable.
  \item Report average clustering and degree assortativity, then visualize one
        complete gold multi-hop path from a development question with node
        types and edge/QA labels. Conclude with a specific explanation of why
        the reconciled projection must not be used by the \PANINI retrieval
        pipeline.
\end{enumerate}

\deliverables At least four plots, two sensitivity tables, and a 200--300 word
cross-dataset interpretation that separates local GSW structure from
reconciliation artifacts.

\grading Connectivity sensitivity (2), degree analysis (2), centrality audit
(2), and clustering/assortativity/path interpretation (2).

\question{4}{Question decomposition and dependency graphs}{10}

The network audit does not yet answer a user's multi-hop question. \RICR first
turns that question into an executable plan. For example, ``Which director
died later?'' may create two independent branches that retrieve each director
and death date, followed by a deterministic comparison. The decomposer emits
text, but the retrieval system needs a graph whose dependencies and operation
types are explicit.

Run the supplied 4B decomposer on all 200 questions using the supplied prompt,
sampling disabled, and the model-config maximum output length. Cache the raw
output before parsing so a parser failure never requires another model run.
Parse every numbered sub-question into a template node. A placeholder such as
\code{<ENTITY\_Q1>} creates an edge from Q1 to the dependent template. Reject
duplicate IDs, missing or future references, cycles, and unresolved
placeholders. Mark retrieval nodes separately from deterministic operations
such as comparison or intersection, then topologically sort the result.

\begin{enumerate}[label=(\alph*)]
  \item Validate references: every placeholder must point to an earlier
        sub-question, the graph must be acyclic, and every retrieval node must
        contain a nonempty natural-language question.
  \item Identify independent linear sequences that may run in parallel and
        record the deterministic operations needed to combine their answers.
  \item On the public reviewed decomposition subset, report valid-plan rate,
        sub-question-count exact match, dependency-edge precision/recall/F1,
        and retrieval-vs-reasoning node accuracy.
  \item Analyze two errors from each dataset. Distinguish malformed output,
        wrong dependency, missing hop, extra hop, and semantically incorrect
        atomic question.
\end{enumerate}

\deliverables \code{decompositions.jsonl}, parser tests, a metric table, and
four concise error analyses.

\grading Inference and caching (2), parser/validation (3), parallel-sequence
handling (2), metrics (2), and error analysis (1).

\question{5}{Sparse retrieval baselines}{8}

The dependency graph now tells the system what atomic question to ask, but not
where to find its answer. Begin with sparse methods because their behavior is
transparent: when a result fails, you can inspect the matched terms and decide
whether the problem is wording, entity description, or graph expansion. These
methods also establish whether later neural components earn their additional
cost.

Give every backend the same return format,
\code{(qa\_uid, score, rank, source)}, and build it first on a tiny hand-written
corpus with one obvious answer. Using only the supplied metadata, implement
and evaluate:

\begin{enumerate}[label=(\alph*)]
  \item TF--IDF retrieval over concatenated QA question and answer text;
  \item BM25 retrieval over the same QA text; and
  \item BM25 entity retrieval over surface name, aliases, roles, and states,
        followed by expansion to QA records attached to each retrieved entity.
\end{enumerate}

Use identical text normalization within a method across datasets. State all
tokenization and BM25 parameters. Evaluate the gold atomic sub-questions from
the development splits at $k\in\{1,5,10,15\}$. Report Recall@$k$ and MRR
overall, by 2Wiki question type, and by MuSiQue hop count.

\deliverables One implementation per method, a complete retrieval table, and
a short explanation of where entity expansion helps or hurts.

\grading Correct TF--IDF/BM25 QA search (3), correct entity expansion (2),
evaluation (2), and interpretation (1).

\question{6}{Dense, hybrid, and paper-style dual retrieval}{12}

Sparse retrieval exposes a central difficulty of multi-hop search: an atomic
question may paraphrase the stored QA text, while an entity mention may be the
best doorway into a relevant local workspace. \PANINI addresses these two
failure modes with two independent routes---dense QA retrieval and sparse
entity retrieval followed by local expansion---and joins them before
reranking. Notice that this expansion uses document-local GSW adjacency from
Question 2, never the analysis-only reconciled graph.

Load the supplied normalized QA embeddings and FAISS index; do not regenerate
corpus embeddings. Use the supplied fixed-query embedding when present. If
answer substitution creates a new query, encode it with the configured query
encoder and cache it by exact normalized query string.

\begin{enumerate}[label=(\alph*)]
  \item Implement dense QA retrieval and verify that a manual inner-product
        calculation agrees with FAISS for at least five queries.
  \item Implement RRF over TF--IDF, BM25-QA, and dense rankings with constant
        60. Deduplicate by stable QA ID.
  \item Implement paper-style dual retrieval: union QA records obtained from
        BM25 entity expansion and dense QA search, deduplicate to a pre-rerank
        pool of at most $M=60$, rerank against the instantiated atomic question,
        and return the top $k=15$.
  \item Compare all methods using the Question 5 protocol. Add mean and
        95th-percentile retrieval latency and the average candidate count.
\end{enumerate}

\deliverables Search code, FAISS consistency test, performance/latency tables,
and the saved top-15 retrieval trace for every development atomic question.

\grading Dense retrieval (3), RRF (2), dual candidate construction (3),
reranking/deduplication (2), and evaluation (2).

\question{7}{Does reranking always help?}{5}

Question 6 deliberately constructs a broad candidate pool: its goal is recall,
so high-ranked items from different routes may be redundant or only
superficially related. The reranker is supposed to repair that ordering using
the instantiated atomic question. It can also overrule a strong retrieval
signal and move the first supporting QA downward. Treat reranking as a
hypothesis to test, not an automatically beneficial step.

Freeze each query's $M=60$ candidate pool so every method sees exactly the same
QA IDs. Compare three scoring rules:
\begin{enumerate}[label=(\roman*)]
  \item reranker probability only;
  \item reciprocal retrieval rank only; and
  \item $0.5$ reranker probability $+\,0.5$ reciprocal retrieval rank.
\end{enumerate}
Use rank reciprocal $1/r$ after assigning each candidate its best pre-rerank
rank. Report Recall@15 and MRR for each dataset. Find one query for which
reranking improves the first relevant rank and one for which it worsens it.
For each, show the top five candidates before and after scoring and explain the
lexical/entity mismatch.

\deliverables One comparison table and two annotated retrieval traces.

\grading Controlled comparison (2), traces (2), and evidence-based explanation
(1).

\question{8}{Implement Reasoning Inference Chain Retrieval}{15}

At this point you have an executable decomposition and a retriever for one
atomic question. The remaining problem is sequential: the answer selected at
one hop changes the text of the next query. Greedily keeping only the
highest-scoring first answer can destroy the correct chain, while retaining
every candidate causes exponential growth. \RICR resolves this tension with a
small beam whose chains carry their own instantiated questions, evidence, and
scores.

Complete \code{prune\_unique\_answers} and
\code{run\_linear\_ricr} in \code{panini\_course/ricr.py}.
Represent a chain as an immutable ordered tuple of evidence records, current
answer, instantiated questions, and hop scores. Implement the search first
with a dictionary-backed toy retriever whose outputs you control, following
this structure:

\begin{lstlisting}[language=Python]
beams = retrieve(first_instantiated_question, k)
beams = prune_unique_answers(beams, B)
for template in remaining_templates:
    expansions = []
    for chain in beams:
        query = substitute(template, chain.current_answer)
        for candidate in retrieve(query, k):
            expansions.append(new_chain(chain, query, candidate))
    beams = prune_unique_answers(expansions, B)
\end{lstlisting}

For each linear sequence:

\begin{enumerate}[label=(\arabic*)]
  \item retrieve $k$ candidates for the first instantiated sub-question;
  \item retain the top $B$ chains with distinct normalized current answers;
  \item substitute each current answer into the next dependent sub-question;
  \item expand all surviving chains by $k$ candidates;
  \item update the geometric-mean chain score in the log domain;
  \item globally prune the $kB$ expansions to $B$ distinct current answers;
  \item repeat until the sequence ends; and
  \item combine parallel sequences by deduplicating evidence using stable QA
        ID while preserving the highest-scoring provenance.
\end{enumerate}

A candidate from an earlier chain may not be mutated in place. Pruning is
global across all $kB$ expansions, not performed separately per parent. Sort by
decreasing geometric-mean score and then stable QA-ID sequence, retaining only
the first chain for each normalized current answer. Ties must use the
deterministic rule from Section 4.2. Return both the deduplicated evidence and
complete trace: instantiated questions, candidates, scores, pruned chains, and
selected QA IDs at every hop. Only after toy cases pass should you connect the
real retriever from Question 7.

\deliverables Working implementation, all supplied \RICR tests passing, at
least three additional tests of your own (including parallel sequences), plus
one human-readable trace from each dataset.

\grading First-hop/pruning behavior (3), multi-hop expansion/substitution (4),
chain scoring (2), parallel evidence merge (2), deterministic trace format
(1), and tests/edge cases (3).

\question{9}{\RICR ablation study}{10}

A working search procedure does not show that its design choices are
necessary. A larger beam can preserve a correct early hypothesis but increases
retrieval calls; unique-answer pruning buys diversity but may discard useful
duplicate evidence; geometric means prevent long chains from winning merely
because they have more terms. Use controlled ablations to determine which
mechanisms produce useful accuracy rather than accidental complexity.

Select a fixed, stratified 20-question subset from each development dataset
and publish the IDs before running experiments. The default is $B=5$, $k=15$,
unique-answer pruning enabled, geometric-mean scoring, and the best frozen
Question 7 retrieval score.

\begin{enumerate}[label=(\alph*)]
  \item Beam width: $B\in\{1,3,5\}$.
  \item Candidates per hop: $k\in\{5,15\}$.
  \item Diversity: unique-answer pruning on versus off.
  \item Chain scoring: geometric mean versus last-hop score.
  \item Retrieval: BM25-QA, dense QA, RRF hybrid, and paper-style dual
        retrieval.
\end{enumerate}

Report supporting-QA recall, complete-chain recovery, answer EM/F1, mean
latency, and mean evidence count. Plot accuracy versus latency and accuracy
versus evidence size. Change only the named factor in each comparison.

\deliverables One ablation table, two trade-off plots, and a 200--300 word
recommendation of the configuration you would deploy.

\grading Experimental control (2), beam/candidate ablations (2),
diversity/scoring ablations (2), retrieval ablation (2), and trade-off analysis
(2).

\question{10}{End-to-end 2Wiki evaluation}{7}

The component experiments now become one complete system in its development
domain. Select the final configuration using only labeled 2Wiki development
questions, record that choice, and freeze it before producing any held-out
prediction. This separation matters: otherwise the final score measures
continued tuning rather than generalization.

Run the frozen system on all 100 2Wiki questions. The answer model may see only
the deduplicated QA evidence returned by \RICR, never source documents,
held-out labels, or reconciled-graph neighbors. Its prompt must request
\code{N/A} when the evidence is insufficient. Save the trace before generating
the answer so retrieval and generation failures can be distinguished.

On the 80 labeled questions, report decomposition validity, supporting-QA
recall, complete-chain recovery, EM, F1, average/95th-percentile latency,
answer-context tokens, and evidence count. Break out EM/F1 by the four 2Wiki
question types. Include one successful and one failed complete trace.

\deliverables \code{results/2wiki\_dev.jsonl},
\code{predictions/2wiki\_heldout.jsonl}, one main table, one type table, and two
traces.

\grading Correct frozen run/output files (2), complete metrics (2),
type-specific analysis (1), trace quality (1), and grounded answer prompting
(1).

\question{11}{MuSiQue transfer and scaling evaluation}{8}

The final test is whether the system learned a reusable procedure or a set of
2Wiki-specific choices. MuSiQue changes both the source distribution and chain
length, including three- and four-hop questions. Treat it as a deployment to a
new collection: package paths may change, but the algorithm and parameters may
not.

Without changing retrieval weights, $M$, $k$, $B$, prompts, normalization, or
score rules, replace the 2Wiki package with the MuSiQue package and run all 100
questions. On the 80 labeled questions:

\begin{enumerate}[label=(\alph*)]
  \item report the same metrics as Question 10, separated into 2-, 3-, and
        4-hop groups;
  \item plot complete-chain recovery and F1 against hop count;
  \item compare 2Wiki and MuSiQue under the same frozen configuration; and
  \item identify whether the largest transfer loss comes from decomposition,
        first-hop retrieval, later-hop substitution/retrieval, or answer
        generation, using counts from saved traces rather than intuition.
\end{enumerate}

\deliverables Submit:
\begin{itemize}
  \item \code{results/musique\_dev.jsonl} and
        \code{predictions/musique\_heldout.jsonl};
  \item a hop-count table and two scaling plots; and
  \item a quantitative transfer-error analysis.
\end{itemize}

\grading Frozen transfer run (2), hop-stratified metrics (2), scaling plots
(1), cross-dataset comparison (1), and traced error attribution (2).

\question{12}{Reproducibility and final submission}{5}

Your final artifact should allow another student to reproduce a result from a
fresh Colab runtime and to inspect where a wrong answer arose. Treat this as a
handoff of the system rather than a notebook snapshot: paths must be relative,
expensive stages restartable, configurations frozen in files, and predictions
separate from gold labels.

Submit one archive or course-repository commit containing:

\begin{enumerate}[label=(\alph*)]
  \item \code{report.pdf}, organized by Questions 1--12 and limited to 12 pages
        excluding references and a one-page appendix;
  \item the completed Colab notebook and any Python modules/tests it imports;
  \item all four development/held-out JSONL files from Questions 10--11;
  \item \code{environment.txt} with Python, CUDA, package, and model revisions;
        and
  \item \code{RUNME.md} with exact commands, seeds, expected runtime, restart
        points, and the final Git commit hash.
\end{enumerate}

The notebook must run from a fresh Colab runtime using relative paths after
setting one \code{PACKAGE\_ROOT}. Do not submit downloaded model weights,
provided embeddings, FAISS indices, or dataset copies.

\grading Report organization/completeness (2), executable code/tests (1),
required result schemas (1), and environment/RUNME reproducibility (1).

\section{Required result schema}

Each prediction file must contain one JSON object per line and preserve the
input order. Extra diagnostic fields are allowed, but the following fields are
required:

\begin{lstlisting}[language=Python]
{
  "dataset": "2wiki" | "musique",
  "split": "development" | "held_out",
  "question_id": "...",
  "question": "...",
  "predicted_decomposition": [...],
  "decomposition_valid": true,
  "retrieval_backend": "dual_hybrid",
  "beam_width": 5,
  "candidates_per_hop": 15,
  "chains": [
    {
      "qa_ids": ["...", "..."],
      "answers": ["...", "..."],
      "hop_scores": [0.91, 0.87],
      "chain_score": 0.8898
    }
  ],
  "evidence_qa_ids": ["...", "..."],
  "predicted_answer": "...",
  "latency_ms": {
    "decomposition": 0.0,
    "retrieval_ricr": 0.0,
    "answer": 0.0
  },
  "answer_context_tokens": 0
}
\end{lstlisting}

Development files may additionally contain gold labels and computed metrics.
Held-out files must not contain gold fields.

\section{Grading summary}

\begin{center}
\begin{tabularx}{0.96\textwidth}{@{}Y r@{}}
\toprule
\textbf{Component} & \textbf{Points}\\
\midrule
Q1. Understand the two data packages & 4\\
Q2. GSW network and entity reconciliation & 8\\
Q3. Structured-memory network analysis & 8\\
Q4. Question decomposition and dependency graphs & 10\\
Q5. Sparse retrieval baselines & 8\\
Q6. Dense, hybrid, and dual retrieval & 12\\
Q7. Reranking analysis & 5\\
Q8. \RICR implementation & 15\\
Q9. \RICR ablations & 10\\
Q10. 2Wiki end-to-end evaluation & 7\\
Q11. MuSiQue transfer and scaling & 8\\
Q12. Reproducibility and submission & 5\\
\midrule
\textbf{Total} & \textbf{100}\\
\bottomrule
\end{tabularx}
\end{center}

\section{General grading and conduct}

\begin{itemize}
  \item A table or plot without axis labels, units, evaluated-question count,
        or a caption receives at most half credit for that item.
  \item Results that cannot be reproduced from the submitted code receive no
        implementation credit, even if the reported values are plausible.
  \item Clearly label any deviation from the required settings. An additional
        experiment does not replace a required one.
  \item You may use standard libraries for graph storage, TF--IDF, BM25, FAISS,
        plotting, and model loading. You may not import an implementation of
        \RICR or copy an answer key.
  \item Follow the course collaboration and academic-integrity policy. Cite
        all external code, prompts, and written sources.
  \item The course late-submission policy applies. If a released artifact is
        corrupted or a Colab dependency becomes unavailable, post a minimal
        reproducer to the course forum before the deadline.
\end{itemize}

\vfill
\begin{center}
\textit{The goal is not merely to obtain a final answer. Your system should
make the network path, evidence, score, and failure mode inspectable.}
\end{center}

\end{document}
